% ---------------------------------------------------------------------------
% Cover letter for Foundations of Physics (Springer Nature).
%
% *** ONE TOGGLE, AND YOU MUST SET IT TRUTHFULLY. ***
%   \prdwithdrawntrue   -> you DID submit to Physical Review D and have withdrawn it
%   \prdwithdrawnfalse  -> you NEVER completed a submission to Physical Review D
% Both compile to a clean letter with no placeholder text. Setting this wrongly
% puts a false statement about submission history in front of an editor, which is
% the single most damaging thing in a cover letter -- so check before compiling.
%
% Compile: tectonic -X compile cover_letter_fop.tex
% ---------------------------------------------------------------------------
\newif\ifprdwithdrawn
\prdwithdrawntrue          % <<<< SET THIS

\documentclass[11pt]{article}
\usepackage[a4paper,margin=1in]{geometry}
\usepackage[T1]{fontenc}
\usepackage{lmodern}
\usepackage{amsmath}
\usepackage{microtype}
\usepackage[hidelinks]{hyperref}
\usepackage{parskip}
\usepackage{enumitem}
\setlist[itemize]{noitemsep,topsep=4pt,leftmargin=1.4em}
\pagestyle{empty}

\begin{document}

\noindent
{\large\textbf{Carl P. Zimmerman}}\\
Briar Creek Tech\\
Charlotte, North Carolina, USA

\vspace{0.6em}
\noindent\hrulefill

\vspace{0.8em}
\noindent 31 July 2026

\vspace{1.2em}
\noindent To the Editors, \textit{Foundations of Physics}

\vspace{0.6em}
\noindent
I submit for your consideration a short research article, \textit{``The MOND acceleration scale is one
half of the dark-energy gravitational rate.''}

\vspace{0.6em}
\noindent\textbf{Why this journal.}
The substance of this paper is conceptual rather than observational, which is why I am offering it to
\textit{Foundations of Physics} rather than to a phenomenology journal. The acceleration scale $a_0$ that
organizes galactic rotation curves has been known since 1983 to lie near $cH_0$, and several closed forms
have been proposed for the coincidence. What I add is not a better fit --- I am explicit that no
improvement in fit is claimed --- but a clarification of what the coincidence actually asserts: written
correctly, it is a statement about \emph{one rational number}, and the apparent geometric structure that
has attached to it dissolves under an elementary reduction. Whether a dimensionless coefficient of order
unity relating a galactic dynamical scale to the cosmological constant can be derived, or must be
accepted as a brute fact, seems to me a foundational question rather than an astronomical one.

\vspace{0.6em}
\noindent\textbf{The result.}
\begin{itemize}
\item $a_0 = \tfrac12 c\sqrt{G\rho_\Lambda} = c^2\sqrt{\Lambda/32\pi} = cH_\Lambda/\sqrt{32\pi/3}$. The
three forms are algebraically identical, and the reduction cancels every factor of $\pi$ and the 3. The
consequence I want to draw attention to is negative: $\sqrt{32\pi/3}$ is \emph{only} a unit conversion
between $H_\Lambda$ and $\sqrt{G\rho_\Lambda}$, so any attempt to explain it as a solid angle, a horizon
count, or a degree-of-freedom tally is reading significance into arithmetic. I make that point partly
because it is an error I made myself.
\item On Planck 2018 inputs the expression gives $9.43\times10^{-11}\,\mathrm{m\,s^{-2}}$, which is
$0.70\%$ from the value obtained by fitting 175 SPARC rotation curves at a stellar mass-to-light ratio
$\Upsilon_{[3.6]} = 0.70$ --- inside the range expected for Spitzer $3.6\,\mu$m populations, so the
agreement does not require an implausible fit.
\item Section~V reports a \emph{failed} attempt to derive the coefficient from ghost-freedom, unitarity
and a holographic degree-of-freedom count. I include the failure because the interesting content of the
paper is the sharpness of the remaining gap, not a claim to have closed it.
\item Section~VI gives a falsifiable consequence: if $a_0\propto\sqrt{\rho_\Lambda}$ then $a_0(z)$ is
exactly constant for $w=-1$ and rises to a maximum before \emph{declining} for the DESI-preferred
$w_0>-1$, $w_a<0$, against the monotonic rise predicted by $a_0\propto cH(z)$. A measured monotonic rise
would exclude the relation.
\end{itemize}

\vspace{0.4em}
\noindent\textbf{What I do not claim, stated here because it bounds the paper.}
The interpolating function I use is \emph{identical} to Eq.~(9) of Milgrom (Phys.\ Lett.\ A \textbf{253},
273, 1999) --- not a variant of it --- and that same paper \emph{derives} $a_0 = 2cH_\Lambda$, a factor
$11.58$ larger. My contribution is therefore a renormalization of his coefficient to match data, not a new
derivation, and Sec.~II says so in those words. Further, the coefficient here is not observationally
distinguishable from Milgrom's later empirical $cH_\Lambda/2\pi$: the two differ by $7.9\%$, which is
$0.018$~dex in $g_{\rm obs}$ against a fit scatter of $0.108$~dex. Section~VII hands over the outstanding
tensions unprompted, including that the cluster-scale discrepancy is $13\%$ \emph{worse} for this lower
$a_0$ than for the conventional value. I would rather a referee found these stated than discovered them.

\vspace{0.6em}
\noindent\textbf{Reproducibility, competing interests, and AI disclosure.}
Every numerical claim is produced by a public script that performs internal consistency checks and exits
with a nonzero status if any fails; the code is archived at
\url{https://doi.org/10.5281/zenodo.21724575}. I am an independent researcher, have no funding to declare,
and have no competing interests. Section~VIII of the manuscript discloses that portions of the analysis
and drafting were carried out with the assistance of a large language model: I directed the work,
specified and reviewed every load-bearing calculation, and take full responsibility including for any
errors. No artificial-intelligence system satisfies the criteria for authorship and none is listed as an
author. Several intermediate claims produced during the work were found to be incorrect and were withdrawn
before submission.

\vspace{0.6em}
\noindent\textbf{Submission history.}
\ifprdwithdrawn
This manuscript was submitted to \textit{Physical Review D} on 31 July 2026 and was withdrawn by me before
review. It is not currently under consideration at any other journal, and it is not part of a joint
submission or a sequence of papers. A preprint is deposited at
\url{https://doi.org/10.5281/zenodo.21724575}; it is not a journal publication. The manuscript is not
posted to arXiv at the time of submission, and I will inform the editorial office if that changes.
\else
This manuscript has not been submitted to any other journal, and it is not currently under consideration
elsewhere. It is not part of a joint submission or a sequence of papers. A preprint is deposited at
\url{https://doi.org/10.5281/zenodo.21724575}; it is not a journal publication. The manuscript is not
posted to arXiv at the time of submission, and I will inform the editorial office if that changes.
\fi

\vspace{0.6em}
\noindent\textbf{Suggested referees.}
The value of this manuscript rests largely on whether its treatment of prior work is judged correct, so I
would particularly welcome a referee familiar with the acceleration-scale literature:
\begin{itemize}
\item Stacy S.\ McGaugh, Case Western Reserve University
\item Beno\^it Famaey, Observatoire astronomique de Strasbourg
\item Federico Lelli, INAF --- Osservatorio Astrofisico di Arcetri
\item Harry Desmond, Institute of Cosmology and Gravitation, University of Portsmouth
\end{itemize}
I have no referees to exclude.

\vspace{0.6em}
\noindent
I am content for the manuscript to be assessed as a short paper.

\vspace{1.2em}
\noindent With thanks for your time,

\vspace{1.4em}
\noindent Carl P. Zimmerman\\
Briar Creek Tech\\
Charlotte, North Carolina

\end{document}
